% Section: Analysis — Why does intensity hurt?
% =====================================================================
% DESCRIPTION
% -----------
% The mechanism story. Without this section the paper is just "we
% tried, it didn't work". With this section it becomes
% "we tried, it didn't work, and here's the architectural reason —
% which the field should know before building variant #47."
%
%   6.1 lambda calibration figure (Figure 2).
%       Source plot: results/lambda_vs_delta.png — copy into
%       paper/figures/. Each point = one (dataset, variant) pair.
%       X = mean learned lambda across seeds. Y = mean test-NDCG@10 Δ
%       vs SASRec across seeds. Reading:
%         - Top-right quadrant: model uses intensity AND gains.
%           Predicted under the hypothesis; observed: empty / sparse.
%         - Bottom-right quadrant: model uses intensity AND loses.
%           Observed: Steam clusters here with lambda 0.55-0.99.
%         - Origin: model opts out (lambda~0) AND recovers SASRec.
%           Observed: a few ML-1M points sit near origin (Val
%           lambda=0.07).
%       The plot's punchline: there is no positive correlation
%       between lambda and Δ. The escape-hatch story FAILS in practice
%       — high lambda does not predict gain (and on Steam, predicts
%       loss).
%
%   6.2 The Steam paradox.
%       Steam has the strongest intensity signal yet the worst
%       outcomes. Two non-mutually-exclusive hypotheses worth probing:
%         (a) Distribution shape: playtime is heavy-tailed
%             (log-normal-ish). After log1p_minmax, a 100h game vs.
%             a 2h game collapses to a narrow band — the model is
%             forced to attend to a near-constant signal that adds
%             variance without information. To diagnose: plot
%             histogram of post-normalisation intensity per dataset
%             (data is in the .inter files). If Steam's distribution
%             is near-constant after norm, the diagnosis lands.
%         (b) Popularity confound: high-playtime games are also
%             popular; item-embedding learning already absorbs that
%             signal, so the extra channel only adds variance.
%             To diagnose: per-dataset Spearman correlation between
%             per-interaction intensity and item popularity rank.
%
%   6.3 Ablation: lambda = 0 vs. lambda = 1 vs. learned.
%       TODO (FLAG IN PAPER): this ablation has NOT been run yet. It
%       is the single most important pre-submission task because it
%       directly tests the "escape-hatch" claim. Expected outcome on
%       Steam: lambda=0 matches SASRec; lambda=1 matches IA-SASRec;
%       learned sits between — confirming that the model could opt
%       out but doesn't, because there is no usable gradient signal
%       pointing it back to lambda=0.
%
%   6.4 Implication for IA-SASRec design.
%       A SINGLE scalar per layer is too coarse:
%       - It cannot down-weight specific intensity values (e.g. mute
%         the 100h tail while keeping the moderate range).
%       - It cannot let some heads use intensity while others ignore
%         it (lambda is shared across heads in our implementation).
%       - It conflates "intensity is informative" with "intensity, in
%         its current normalised form, is informative" — confounding
%         the architectural claim with the preprocessing choice.
% =====================================================================
\section{Analysis}
TODO: 6.1 lambda-vs-Delta figure + reading, 6.2 Steam paradox probes,
6.3 lambda ablation (mark as REQUIRED before submission), 6.4
architectural implications.
